\chapter{Definición de criterios de evaluación y recolección de datos}
\label{ch:recoleccion}

\section{Recolección de Instancias}
\label{sec:recoleccion_instancias}

A diferencia de áreas como la programación lineal entera, donde repositorios como MIPLIB
concentran miles de instancias bajo un formato y criterio de calidad unificados \cite{gleixnerMIPLIB2017Datadriven2021}, 
la programación semidefinida carece de un repositorio análogo de referencia única. Las instancias de \acrshort{sdp} 
se encuentran dispersas entre colecciones mantenidas de forma independiente, con formatos heterogéneos, escalas 
distintas y orígenes aplicados muy variados. Esta fragmentación obligó a extraer instancias de múltiples fuentes y 
normalizarlas a un formato común antes de poder realizar cualquier análisis comparativo.

La calidad de un análisis del espacio de instancias depende directamente de que estas cubran de forma amplia y
representativa el espacio de problemas posibles. Una base de instancias sesgada
hacia una sola fuente o tipo produciría un espacio de instancias con regiones artificialmente vacías que reflejarían
la ausencia de diversidad experimental, no una característica genuina del problema \cite{smith-milesInstanceSpaceAnalysis2023}.

Se seleccionaron cuatro bibliotecas de benchmarks y librerias existentes con criterios de diversidad estructural y aplicada:

\textbf{SDPLIB 1.2}\cite{borchersSDPLIB12Library1999}. La biblioteca de referencia histórica del área, con
instancias provenientes de control óptimo, diseño estructural, localización de antenas y relajaciones semidefinidas
de problemas combinatoriales. Si bien intenta ser la librería de referencia estándar para problemas de \acrshort{sdp}
de igual manera que MIPLIB lo es para los problemas de Programación Entera Mixta, esta solo cuenta con 92 instancias en la actualidad,
lo cual limita su variedad.

Su amplio uso en la literatura facilita la comparación con trabajos previos.

\textbf{Benchmarks de Hans Mittelmann}\cite{PlatoasueduFtpSparse_sdphtml}. Colección de instancias de
mediana y gran escala mantenida por Hans D. Mittelmann, investigador reconocido por sus estudios comparativos
de software de optimización numérica. Sus benchmarks publican instancias recopiladas desde distintas fuentes
académicas y aplicaciones de optimización, incluyendo los problemas \texttt{kocvara\_pennon}, provenientes
de optimización estructural con variables de forma. Esta fuente complementa SDPLIB con instancias de mayor
dimensión que estresan el comportamiento de los solvers en regímenes de memoria y tiempo.

\textbf{DIMACS Library}\cite{DIMACSLibraryMixeda}. Colección asociada al DIMACS
(\textit{Center for Discrete Mathematics and Theoretical Computer Science}), centro de investigación
dedicado a matemática discreta y ciencias de la computación teórica. En particular, estas instancias
provienen del \textit{7th DIMACS Implementation Challenge}, orientado a evaluar el avance de los solvers
de propósito general mediante problemas de programación semidefinida-cuadrática-lineal mixta. Esta fuente
aporta instancias con mezcla de variables continuas y semidefinidas.

\textbf{Biq Mac Library}\cite{BiqMacLibrarya}. Biblioteca asociada a Biq Mac, un software especializado
para resolver problemas binarios cuadráticos y de corte máximo mediante técnicas de optimización semidefinida
y métodos de ramificación y acotamiento. Sus instancias corresponden a problemas de corte máximo
(\textit{max-cut}) formulados como relajaciones \acrshort{sdp}, con estructura de grafos densos.
Constituyen una clase estructuralmente distinta: una sola restricción por arista, bloque único de gran
dimensión y densidad alta. Esto introduce una región del espacio de instancias que las otras tres fuentes
no cubren.

En conjunto, las cuatro fuentes abarcan orígenes aplicados heterogéneos ---control, geometría, combinatoria, diseño
estructural--- y un rango amplio de valores para las features, lo que favorece la separabilidad de regiones
en el espacio de instancias.

Cabe notar que no todas las instancias de estas fuentes corresponden a programas semidefinidos puros en su formulación
original: en varios casos provienen de problemas de optimización más generales ---mixtos, cuadráticos o combinatoriales---
cuya relajación semidefinida constituye la instancia efectivamente analizada. Esta práctica es estándar en el área y no
compromete la validez del análisis, dado que las características se extraen directamente de la formulación SDP
resultante, independientemente del origen del problema.

\section{Métrica de Rendimiento y Criterio de Buen Desempeño}
\label{sec:metrica}

El tiempo de \acrshort{cpu} fue seleccionado como métrica principal de rendimiento por ser la medida más directamente comparable entre solvers de
distintas familias algorítmicas e implementaciones. Las métricas alternativas, como el número de iteraciones, son incomparables entre métodos de
punto interior y métodos de Lagrangiano aumentado dado que cada iteración tiene un costo muy diferente.

El etiquetado binario de rendimiento (bueno/malo) sigue el criterio de umbral relativo de \acrshort{isa}:
un solver se considera \textit{bueno} en una instancia si su tiempo de \acrshort{cpu} no supera el $55\%$ del tiempo del solver más rápido en esa
instancia ($\beta = 0.55$). Este umbral, recomendado por la metodología, tolera una penalización moderada respecto al óptimo sin considerar
competitivos a solvers que resulten significativamente más lentos. Las instancias para las que ningún solver converge dentro del límite de
$1200$ segundos se excluyen del análisis de rendimiento.